\section{Discussion and Limitations}\label{sec:disc}

\noindent\textbf{What the Guarantees Cover.} The temporal-consistency claim of \sname\ rests on P2, a deterministic and pathwise lower bound on the interval between discretionary switches that needs only the range bound of the normalized cost and the declared non-vacuity condition \eqref{eq:p2nonvac}; P1(ii) sharpens it to exponential rarity only under an i.i.d.\ hypothesis that stationarity alone does not supply (\cref{rem:dep}). Because P2 is satisfied by a policy that never switches, it is paired with the conditional response-time bound of \cref{prop:p2live}, and the ablations show both sides of that pairing in data: the leaky accumulator suppresses argmin oscillation in \texttt{near\_tie}, while at the same default knobs the late reversal in \texttt{mid\_reversal} is not tracked within the log (an observed outcome that the response bound does not predict, and that only an independent opportunity oracle could score as a miss) and the fast-saturating progress input of \cref{fig:rhosweep} triggers clip-induced lockout (\cref{subsec:ablation}). P3 and P5 are branch-priority and guard contracts, not physical collision-freedom or termination guarantees (\cref{prop:p3,prop:p5}). All of these apply standard dwell-time and CUSUM/queueing arguments (\cref{sec:related,subsec:props}); their value lies in stating, for this rule, the assumptions under which each holds, not in new mathematics.

\noindent\textbf{Cost Analysis.} The historical decision-core timing of \cref{tab:latency} leaves the per-cycle decision computation well inside a 20~Hz period at the default $K_{\rm cap}=3$, but it is a finite-batch measurement of \code{DecisionCore::step} only; the candidate-trajectory adapter, live costmap access, and the end-to-end ROS~2 pipeline are not timed (\cref{subsec:latency}).

\noindent\textbf{Limitations.} (i) The planar (2D) assumption means multi-level or sloped environments are not handled. (ii) Constant-velocity Kalman prediction is strong in the short term but has limits under abrupt intent changes or multi-modal crowds (a learning-based predictor can be substituted as an interchangeable input, \cref{sec:related}). (iii) Front/rear passing is handled only as cost and is not separated into distinct \Class{es}. (iv) The empirical scope of this paper is limited to the decision layer. Behavioral consistency, the decision-stream observer statistic, and per-cycle latency were \textbf{measured} with an offline harness driving the actual decision core (\cref{subsec:main,subsec:ablation,subsec:latency}), but performance metrics requiring physical simulation (success rate, collision, minimum social distance), the external baselines, and the user study have not yet been measured in this evidence set; their protocols have been fixed in advance in the Planned Evaluation (\cref{app:planned}), with per-case values still to be bound before any run. Meanwhile, the ablations showed that the leaky accumulator is the primary oscillation-suppression mechanism (removing $-$hysteresis or $-$hardening alone produced no regression in this regime), and revealed an offline transition-blocking interval for the scripted $\rho$ drive rate ($v_{\text{lat}} \in (0.20, 0.35)$ m/s), not a physical freezing threshold. (v) \textbf{No consistency-weight, opinion-dynamics or closed-loop comparison.} The comparators are internal (immediate argmin, simple dwell, knob ablations and a synthetic correspondence-loss stress comparator). The most direct missing baseline is a previous-class consistency weight~\citep{degroot2024topology,martinezbaselga2026shine} applied to exactly the same candidates and cost streams; opinion-dynamics passing~\citep{cathcart2023opinion} and dwell-plus-relative-improvement selection~\citep{li2026bigcbf} are further baselines. None of these was run and no closed-loop simulation was performed, so the offline results show only that the implemented rule filters scripted, noisy decision-label streams; whether it improves on these mechanisms, or on navigation outcomes, is untested. These comparisons are part of the planned evaluation (\cref{app:planned}). (vi) \textbf{Strength of the guarantees.} At the default knobs the worst-case separation guaranteed by P2 is short (\cref{rem:p2nonvac}), P1 relies on an i.i.d.\ hypothesis that temporally correlated social costs need not satisfy (\cref{rem:iid}), and a policy that never switches satisfies P2; the guarantees are contracts about the decision rule, not evidence of legible or effective navigation. The current contribution is therefore the conditional decision-rule design and formal properties together with \textbf{offline empirical validation of decision-layer consistency}; physical performance, external comparison, and human-data validation are to be filled in subsequently according to the protocol of \cref{app:planned}.

\noindent\textbf{Generalization limits (P4, actuation, sensing).} P4's guarantee of maneuver completion depends on $v_{\text{lat,min}} > 0$ in \cref{asm:scope}(d), which presupposes the following. (i) \textbf{Actuation type.} For nonholonomic drives (differential, Ackermann), cross-track progress is realized only as the product of forward speed and steering, so $v_{\text{lat,min}} > 0$ requires \textbf{positive forward speed and forward clearance}. That is, if the robot stops or the front is blocked, the lateral progress rate becomes 0, breaking P4's premise, and the safety layer (P3, P5) takes over. For holonomic (omnidirectional) drives, direct lateral motion is possible even while stationary, relaxing this premise. (ii) \textbf{Footprint.} The corridor-width and pruning analysis in this paper assumes a circular footprint approximation; a rectangular footprint (an elongated body) has an occupied area that varies with orientation during rotation, which can make corridor-availability determination and the realization of $v_{\text{lat,min}}$ more conservative. (iii) \textbf{Dependence on forward speed.} As noted in (i), $v_{\text{lat,min}} > 0$ depends on forward speed, so in congested environments with frequent low-speed or stopped intervals, the upper bound $K$ on maneuver-completion time increases. (iv) \textbf{Sensor modality.} This design receives person tracking, F-formation, and vulnerable-pedestrian attributes from external modules; 2D-laser-based tracking has weak heading and attribute estimation, whereas depth/RGB-D-based tracking is richer. The quality of $\sigma_{\text{front}}/\sigma_{\text{rear}}$ regime determination and of group/attribute recognition therefore depends on sensor modality, and the guarantees of this paper's decision layer hold on the premise that these inputs are of adequate quality.

\noindent\textbf{Ethical limitations.} The vulnerable-pedestrian handling of \cref{subsec:groups} relies on external attribute classification, which carries risks of bias, privacy, and stigma. This paper specifies only a policy of choosing conservative defaults under uncertainty and does not claim the fairness of the classifier itself as a contribution, and it explicitly notes that a separate ethics and privacy review is needed for real deployment.

\noindent\textbf{Future work.} Carrying out the Planned Evaluation (\cref{app:planned}) --- physical-simulation performance metrics, independently tuned comparison against the external baselines, and validating time-to-legible against human data via the user study --- is the top priority; beyond that, we propose extending front/rear passing into explicit \Class{es} via spatiotemporal $(x,y,t)$ homotopy, and combining the approach with learning-based prediction.
